% =========================================================
% IEEE Two-Column Conference Paper
% Title: From Hand-Crafted Heuristics to Learned Representations:
%        A Survey of Object Detection Paradigm Shifts
% Format: IEEEtran, two-column conference style
% =========================================================

\documentclass[conference]{IEEEtran}

\usepackage{cite}
\usepackage{amsmath,amssymb,amsfonts}
\usepackage{graphicx}
\usepackage{textcomp}
\usepackage{xcolor}
\usepackage{hyperref}
\usepackage{booktabs}

\hypersetup{
  colorlinks=true,
  linkcolor=blue,
  citecolor=blue,
  urlcolor=blue
}

\begin{document}

% ---------------------------------------------------------
% Title and Author Block
% ---------------------------------------------------------
\title{From Hand-Crafted Heuristics to Learned Representations:\\
A Survey of Object Detection Paradigm Shifts}

\author{
  \IEEEauthorblockN{Author Name}
  \IEEEauthorblockA{
    \textit{Department of Computer Science}\\
    \textit{University Name}\\
    City, Country\\
    author@university.edu
  }
}

\maketitle

% ---------------------------------------------------------
% Abstract
% ---------------------------------------------------------
\begin{abstract}
Automated object detection---the task of simultaneously localizing
and classifying instances of semantic categories within an image---has
undergone a fundamental transformation over the past two decades.
Prior to 2012, practitioners relied on carefully constructed feature
descriptors and statistical classifiers that encoded domain knowledge
about edges, textures, and spatial arrangements.
The widespread adoption of deep convolutional neural networks after
2012 displaced this paradigm almost entirely, replacing hand-crafted
pipelines with hierarchical feature learning end-to-end.
This survey traces that transition in three phases:
(i) the classical era of Viola-Jones, HOG+SVM, and deformable
part-based models; (ii) the first wave of deep-learning detectors
anchored by region proposals; and (iii) the single-pass architectures
that brought real-time performance to competitive accuracy.
We analyze the technical motivations behind each phase shift,
identify the bottlenecks that necessitated the next advance, and
discuss the benchmark trajectories that quantify the gains.
\end{abstract}

\begin{IEEEkeywords}
object detection, convolutional neural networks, Viola-Jones,
histogram of oriented gradients, deformable part models,
deep learning, PASCAL VOC, image recognition
\end{IEEEkeywords}

% =========================================================
% PAGE 1 --- Introduction
% =========================================================
\section{Introduction}
\label{sec:intro}

Object detection is among the most demanding problems in computer
vision because it couples two distinct challenges: determining
\emph{where} objects appear in an image (localization) and deciding
\emph{what} those objects are (classification).
A detector must generalize across wide intra-class variation---the
same category can span orders of magnitude in scale, exhibit
dramatic pose change, and appear under partial occlusion or
adverse lighting.

Early attempts at solving detection focused on specific,
well-constrained categories such as human faces and pedestrians,
for which it was feasible to enumerate the relevant visual cues by
hand.
Researchers encoded these cues as feature vectors computed over
image patches, then trained shallow classifiers to separate positive
from negative examples in the resulting feature space.
This strategy, though effective for narrow domains, struggled as
the breadth of target categories expanded.

The introduction of standardized benchmarks transformed the field
into an empirical discipline.
The PASCAL Visual Object Classes challenge \cite{everingham2010pascal},
launched in 2005, defined a canonical evaluation protocol---mean
Average Precision (mAP) across twenty categories---that allowed
direct, reproducible comparison of detection systems.
Year-on-year incremental improvements under PASCAL VOC revealed
both the remarkable ingenuity of classical engineering and its
structural ceiling.

That ceiling was shattered in 2012 when Krizhevsky \emph{et al.}
demonstrated that a deep convolutional network trained on ImageNet
could learn feature representations substantially more powerful than
anything hand-designed \cite{krizhevsky2012imagenet}.
Within two years, Girshick \emph{et al.} adapted this insight to
detection through the Region-based CNN (R-CNN) framework
\cite{girshick2014rich}, establishing deep features as the new
backbone of high-accuracy detection.
Subsequent architectures---Faster R-CNN \cite{ren2015faster},
SSD \cite{liu2016ssd}, and YOLO \cite{redmon2016you}---refined
the efficiency and end-to-end trainability of these deep pipelines,
culminating in detectors that operate in real time at accuracies
that would have seemed unreachable in the pre-deep-learning era.

The remainder of this paper is organized as follows.
Section~\ref{sec:traditional} examines the classical detection
methods and their underlying assumptions.
Section~\ref{sec:deep} surveys the deep-learning-based detectors
that replaced them.
Section~\ref{sec:analysis} offers a comparative analysis of the
two eras, and Section~\ref{sec:conclusion} draws broader conclusions
about the nature of representation learning in computer vision.

% =========================================================
% PAGE 2 --- Foundations of Computer Vision (Traditional Methods)
% =========================================================
\section{Foundations of Computer Vision:\\Traditional Methods}
\label{sec:traditional}

Before gradient-based feature learning became viable, detection
systems were constructed around two separable design decisions:
how to represent a candidate region as a numeric vector, and
how to discriminate positive from negative vectors.
Three methodological families---each addressing a distinct category
and scale of problem---defined the pre-2012 landscape.

% ---------------------------------------------------------
\subsection{The Viola-Jones Detection Framework}
\label{subsec:vj}

The framework introduced by Viola and Jones
\cite{viola2001rapid,viola2004robust} represented a watershed
moment for face detection.
Its central contribution was not a single algorithmic idea but
an integrated architecture that made real-time performance
achievable on the hardware of 2001.

\subsubsection*{Haar-like Rectangular Features}

The feature vocabulary of the framework consists of
\emph{Haar-like features}---difference-of-rectangle computations
evaluated at specific positions and scales within a detection
window \cite{lienhart2002extended}.
A Haar-like feature $h_{j}$ anchored at position $(x, y)$ with
width $w$ and height $h$ is computed as:
\begin{equation}
  h_{j}(x,y,w,h) = \sum_{(i,k)\in R_{+}(x,y,w,h)} I(i,k)
                 - \sum_{(i,k)\in R_{-}(x,y,w,h)} I(i,k),
  \label{eq:haar}
\end{equation}
where $R_{+}(x,y,w,h)$ and $R_{-}(x,y,w,h)$ denote the white and
grey rectangular sub-regions determined by the anchor parameters,
and $I(i,k)$ is the pixel intensity at location $(i,k)$.
Enumerating all possible positions, scales, and orientations
yields over 160,000 distinct features for a $24 \times 24$ window.

\subsubsection*{Integral Image Acceleration}

Evaluating even a small fraction of those features naively would
be prohibitively expensive.
The \emph{integral image} (also called a summed-area table)
transforms the computation so that any rectangular sum requires
exactly four memory lookups, regardless of rectangle size.
Formally, the integral image $\text{ii}(x,y)$ at position $(x,y)$
accumulates all pixel values above and to the left:
\begin{equation}
  \text{ii}(x,y) = \sum_{x' \leq x,\; y' \leq y} I(x', y').
  \label{eq:ii}
\end{equation}
This constant-time summation is what allows the framework to
evaluate thousands of features per window with low latency.

\subsubsection*{AdaBoost Feature Selection and Cascade}

Because evaluating all features at every scale and position
remains costly, Viola and Jones apply AdaBoost
\cite{freund1997decision} to select a small, highly discriminative
subset, forming a \emph{strong classifier} from a weighted
combination of weak decision stumps.
These strong classifiers are then arranged in an \emph{attentional
cascade}: each stage rejects clear negatives early, so that the
computationally expensive later stages are invoked only for
promising candidate regions.
In practice, over 99.9\% of sub-windows are discarded by the
early stages, yielding detection speeds that matched or exceeded
real-time requirements at the time of publication.

% ---------------------------------------------------------
\subsection{Histograms of Oriented Gradients with SVM}
\label{subsec:hog}

Pedestrian detection introduced challenges absent from face
detection: highly variable clothing, articulated limb positions,
and a much wider range of aspect ratios.
Dalal and Triggs addressed these challenges through the
\emph{Histogram of Oriented Gradients} (HOG) descriptor
\cite{dalal2005histograms}, which encodes the local distribution
of gradient directions rather than absolute pixel intensities.

\subsubsection*{Gradient Computation and Spatial Binning}

Given an image patch, the spatial gradient magnitude
$|\nabla I(x,y)|$ and orientation $\theta(x,y)$ are computed at
every pixel using finite-difference filters.
The detection window is divided into a grid of small spatial
\emph{cells} (typically $8 \times 8$ pixels), and within each
cell the gradient orientations are soft-binned into a histogram
of $K$ orientation channels (commonly $K=9$, spanning $0^\circ$
to $180^\circ$ unsigned).
Adjacent cells are grouped into larger \emph{blocks}
($2 \times 2$ cells), and each block histogram is
$\ell_2$-normalized to suppress the effect of local illumination
changes.
The concatenation of all normalized block histograms forms the
final HOG descriptor.

\subsubsection*{Linear SVM Classification}

The high-dimensional HOG vector is fed to a linear Support Vector
Machine trained to produce a binary pedestrian/non-pedestrian
decision.
The SVM maximises the margin between the hyperplane separating
positive training windows from hard-negative examples mined from
background imagery.
Dalal and Triggs reported a miss rate below 10\% at $10^{-4}$
false positives per window on the INRIA dataset, far surpassing
competing descriptors of the era.
The coupling of HOG with a linear SVM became the canonical
pedestrian detection pipeline for the following several years.

% ---------------------------------------------------------
\subsection{Deformable Part-Based Models}
\label{subsec:dpm}

Even with rich gradient features, single rigid templates
proved insufficient for categories exhibiting structural
deformation---animals in motion, bicycles at varying viewing
angles, or people in non-canonical poses.
Felzenszwalb \emph{et al.} addressed this through the
\emph{Deformable Part-based Model} (DPM)
\cite{felzenszwalb2010object,felzenszwalb2008discriminatively},
which represents an object as a coarse root filter combined with
a set of higher-resolution part filters positioned at learned
offsets from the root.

\subsubsection*{Mixture Models and Latent SVM}

Each object category is described by a mixture of $K$ component
models, where each component captures a specific viewpoint or
configuration (e.g., a frontal versus profile view of a car).
The part filters and their spatial deformation costs are jointly
learned using a \emph{Latent SVM} (LSVM), which treats part
positions as latent variables.
The LSVM objective alternates between fixing latent assignments
to compute supergradients and re-estimating classifier weights,
converging to a locally optimal set of filter templates and
deformation penalties.

\subsubsection*{Scoring and Detection}

At test time, the score of a candidate bounding box is computed
as the root filter response plus the maximum over all valid
part placements of the sum of part filter responses minus
deformation penalties, with $n$ denoting the total number of
part filters:
\begin{equation}
  S(\mathbf{p}) = F_0 \cdot \phi(I, p_0)
    + \sum_{i=1}^{n}\bigl[F_i \cdot \phi(I,p_i)
    - d_i \cdot \psi(p_0,p_i)\bigr],
  \label{eq:dpm}
\end{equation}
where $F_0$ is the root filter, $F_i$ are the $n$ part filters,
$\phi(I, p)$ denotes the HOG feature vector at position $p$,
$d_i$ encodes the quadratic deformation cost, and
$\psi(p_0, p_i)$ measures spatial displacement.
DPM dominated the PASCAL VOC leaderboard from 2008 through
roughly 2012, when deep features began to supersede it.

% ---------------------------------------------------------
\subsection{Limitations of Hand-Crafted Feature Pipelines}
\label{subsec:limits}

Despite representing the best engineering efforts of their era,
the classical methods described above share a fundamental
architectural constraint: the feature representation is
predetermined by human designers rather than inferred from
data.

Haar-like features express pairwise intensity contrasts but
carry no explicit model of curvature, symmetry, or high-order
texture.
HOG captures local gradient statistics yet aggregates them
over fixed spatial grids that cannot adapt to object boundaries.
DPM part filters, though learned from examples, are still
linear filters applied to hand-crafted HOG maps, so their
expressive capacity is bounded by the underlying descriptor.

Collectively, these representations enforce an \emph{inductive
bias} that aligns well with the specific categories for which
they were designed but generalises poorly beyond them.
The result was a hard ceiling in achievable detection accuracy:
systems trained on PASCAL VOC plateaued near 40--45\% mAP
despite years of incremental refinement, because no amount of
classifier tuning could compensate for the descriptive gaps
inherent in manually defined feature spaces.
The inability to compose low-level filters into abstract
semantic representations---capturing, say, the concept of a
partially occluded bicycle wheel at an arbitrary viewing
angle---limited detection precision in precisely the situations
that matter most in real-world deployment.
This limitation set the stage for a fundamentally different
approach: end-to-end learning of hierarchical representations
directly from labelled image data, which is the subject of
the following section.

% =========================================================
% References
% =========================================================
\bibliographystyle{IEEEtran}
\bibliography{references}

\end{document}
